在 \(L^2([0,1])\) 中，用常数函数逼近 \(f(x)=x\) 的最佳均方近似是 \(c=\int_0^1 x\,dx=\frac12\)，误差 \(\|f-c\|_2^2=\frac1{12}\)。这是正交投影的最简实例：在常数子空间上的投影就是期望。

前面的单元把函数看作被积分、被逼近的对象。换一个眼光：把每个函数本身当成空间中的一个点，问两点之间有多远、一串点是否收敛、极限是否还在这个空间里。这一步把分析从``研究函数''升级为``研究函数组成的空间''。

本单元回答四个问题：

\begin{enumerate}
\tightlist
\item
  \(L^p\) 范数把积分变成距离：不同 \(p\) 各强调误差的哪一面，H\"older 与 Minkowski 不等式提供哪些基本控制。
\item
  完备性让 \(L^p\) 成为 Banach 空间：Cauchy 函数列的极限为什么仍在 \(L^p\) 中，这个性质支撑哪些构造。
\item
  \(L^2\) 的几何让投影与正交可用：内积、正交投影、Riesz 表示与 Fourier 级数怎样放进同一框架，RKHS 的求值泛函为何连续。
\item
  条件期望与线性回归是两种投影：最佳均方预测、最佳线性预测各自是什么，Gaussian 情形为何重合。
\end{enumerate}

选范数就是选``什么算接近''。Part 03 的\hyperref[sec:12-Real-Analysis-Support/03-Topology-Basics/01]{``距离怎样产生开集与闭集''}在有限维说这件事，本单元在函数空间把它落地：\(L^2\) 小不控制最坏点，sup 小却控制处处；下一单元的条件期望之所以是``最佳预测''，正是因为 \(L^2\) 几何中的正交投影定理。

\subsection{一个具体算例}

\(L^p([0,1])\) 上取 \(f_n(x)=x^n\)：\(\|f_n\|_p=(\int_0^1 x^{np}\,dx)^{1/p}=(np+1)^{-1/p}\to 0\)（\(p<\infty\) 时），函数列依 \(L^p\) 范数收敛到 0。但 \(\|f_n\|_\infty=1\) 对所有 \(n\)——一致范数下不收敛。一个函数列可以在 \(L^p\) 中消失却在最大值意义下不变。

\(L^2([0,2\pi])\) 中，\(\{1,\cos nx,\sin nx\}_{n=1}^\infty\) 是正交基。\(f(x)=x\) 的 Fourier 系数为 \(a_0=\pi\), \(a_n=0\), \(b_n=-2/n\cdot(-1)^n\)，截断到 \(N=5\) 的投影已能捕获主要形状。

\subsection{怎么读这一章}

先用 \(L^p\) 范数把函数差异变成距离，再看完备性为何使极限仍是可用函数。只有到 \(L^2\) 时才出现真正的内积、正交和投影；最后一篇用条件期望与线性回归对比两种投影空间。

\subsection{本章篇目}

以下篇目按概念依赖排列；若从具体问题进入，也可以先打开最接近当前困惑的一篇，再沿篇内链接回补。

\begin{itemize}
\tightlist
\item \hyperref[sec:12-Real-Analysis-Support/07-Lp-Spaces-and-Function-Geometry/01]{``\(L^p\) 范数把积分变成距离''}；
\item \hyperref[sec:12-Real-Analysis-Support/07-Lp-Spaces-and-Function-Geometry/02]{``完备性让 \(L^p\) 成为 Banach 空间''}；
\item \hyperref[sec:12-Real-Analysis-Support/07-Lp-Spaces-and-Function-Geometry/03]{``\(L^2\) 的几何让投影与正交可用''}；
\item \hyperref[sec:12-Real-Analysis-Support/07-Lp-Spaces-and-Function-Geometry/04]{``条件期望与线性回归是两种投影''}。
\end{itemize}
